\chapter{Hashtries}\label{chap:hashtries}

Hashtries (Hash Tries) are a trie-based data structure that uses the bits of the key's hash to navigate the tree. They offer a good balance between the speed of hashtables and the structural properties of trees.

\section{Design rationale}

Our hashtrie implementation uses a recursive structure. Each node in the trie represents a prefix of the hash. We use a branching factor of $N=4$ (processing 2 bits of the hash at each level) or $N=16$ (4 bits).

Collisions (different keys with the same hash) are handled by storing a list of elements at the leaf nodes. Since the tree depth is limited by the number of bits in the hash (e.g., 64 bits), the worst-case lookup time is effectively constant, provided the hash function distributes keys well.

Hashtries do not require resizing in the same way as hashtables, as they grow incrementally. This avoids the latency spikes associated with rehashing large tables.
